\section{The new continuation shortcut and a direct iPRA theorem}
\label{sec:shortcut}
Let $L$ be either new loop-reassociated parent, before the three shortcuts.
Let $F$ be its accelerated version, the delivered machine. All labels in
this section are native labels. Write configurations as $(q,h,T)$ and put
$T_j=T(h+j)$. The following finite invariant suffices:
\begin{align*}
q=101&\Longrightarrow T_{-1}=0,\\
q=11&\Longrightarrow (T_{-1}=0\lor T_1=0),\\
q=26\ \land\ T_0=1&\Longrightarrow T_1=T_2=0,\\
q=190&\Longrightarrow T_{-1}=1.
\end{align*}
The initial control is 74, so these implications initially hold for every
complete tape. The table checker verifies that the only predecessor of
101 is $96/0\mapsto0R101$, every transition into 11 writes zero, the only
predecessors of 26 are $11/0\mapsto0L26$ and $101/0\mapsto0L26$, and the
only predecessor of 190 is $139/1\mapsto1R190$.

\begin{lemma}[Incoming-symbol invariant]
Every run of $L$ beginning in control 74 satisfies the four implications.
\end{lemma}
\begin{proof}
The first, second and fourth assertions follow directly from the cell just
written by an incoming transition. For the third, entering 26 writes zero
immediately to the new head's right. An entry from 101 scans zero at the
new head, by the first assertion, so its read-one case cannot occur. An
entry from 11 that scans one makes the old left neighbor one. The second
assertion then forces the old right neighbor to be zero: this is the new
cell at offset two. Induction on the number of transitions completes the
proof. The assertion is a Boolean combination of finitely many tape-bit
functions and hence is quantifier free in the primitive-recursive language.
\end{proof}

The two older shortcuts from the incoming acceleration packet are
\begin{align*}
(26,h,T)&\xrightarrow{1R42}(42,h+1,T)
 \xrightarrow{0R13}(13,h+2,T)\xrightarrow{0L91}(91,h+1,T),\\
(101,h,T)&\xrightarrow{0L26}(26,h-1,T)
 \xrightarrow{0R194}(194,h,T)\xrightarrow{0R17}(17,h+1,T).
\end{align*}
The guards are respectively $T_0=1,T_1=T_2=0$ and $T_{-1}=T_0=0$.
They give the replacements $26/1\mapsto1R91$ and
$101/0\mapsto0R17$.

\textbf{The additional shortcut found here} uses $T_{-1}=T_0=1$:
\[
(190,h,T)\xrightarrow{1L73}(73,h-1,T)
 \xrightarrow{1R83}(83,h,T)\xrightarrow{0L42}(42,h-1,T[h:=0]).
\]
It is exactly the single transition $190/1\mapsto0L42$. This change does
not follow from state renaming: it erases a bit that the previous first
transition preserved, and removes two physical transitions on every
applicable visit. The entire exterior tape is untouched. No intermediate
control displayed here is the halt control 254. Native control 194 in the
second identity is a working control, not the terminal label used by dense
numbering.

The finite checker verifies all structural premises and all twenty
length-five tape contexts satisfying the three guards: four for 26,
eight for 101, and eight for 190. Each local equality compares the full
five-cell tape, the head, and the control; the proof above explains why
arbitrary exterior cells are irrelevant. The invariant, rather than a
finite run sample, licenses their use at arbitrary times.

\begin{theorem}[Exact expanded clock]
Totalize each step function by making the halting configuration absorbing.
For a common initial configuration with control 74 there is a primitive-
recursive clock $\sigma$ satisfying
\[
C^F_t=C^L_{\sigma(t)},\qquad t\le\sigma(t)\le3t.
\]
In particular, with $H_M(t)$ meaning halted by time $t$, iPRA proves
\[
H_L(t)\longrightarrow H_F(t),\qquad
H_F(t)\longrightarrow H_L(3t).
\]
\end{theorem}
\begin{proof}
Let $k(C)=3$ at the three changed source/symbol pairs and $k(C)=1$
elsewhere, including the absorbing halt. Define
$\sigma(0)=0$ and $\sigma(t+1)=\sigma(t)+k(C^F_t)$.
Induct on $t$ using the preceding invariant and exact local identities.
Each new step equals the specified positive number of old steps; the old
invariant remains true at the next matching configuration. The clock bounds
follow from $1\le k\le3$. All functions are primitive recursive and the
induction predicates are quantifier free.

If $L$ has halted by $t$, its state at $\sigma(t)\ge t$ is still halted,
so $F$ is halted at $t$. If $F$ is halted at $t$, then $L$ is halted at
$\sigma(t)\le3t$ and hence at $3t$. This argument uses absorbing semantics
only in the arithmetic representation; it adds no working state to the
literal machine. It applies to arbitrary finite-support initial tapes,
not only to the blank initialization used for the consistency machines.
\end{proof}

\section{Complete physical verification and the relative source theorem}
\label{sec:relative}
The loop reassociation changes addresses, so the two literal machines are
compared through their common resolved primitive program. This section
spells out the weak-arithmetic argument that connects the checked finite
contracts to the all-time theorem. It does not infer nonhalting from a
bounded simulation.

\subsection{Arithmetic representation and the checker rules}
A tape zipper is a pair $(l,r)$ of natural numbers. The low bit of $r$ is
at the head and the low bit of $l$ is immediately to its left. Writing $w$
and moving right gives
\[
(l,r)\longmapsto(2l+w,\lfloor r/2\rfloor);
\]
writing $w$ and moving left gives
\[
(l,r)\longmapsto
 (\lfloor l/2\rfloor,\ 4\lfloor r/2\rfloor+2w+(l\bmod2)).
\]
These formulae encode equality of the complete finite-support tape,
relative to the head, including its two blank exteriors. An additional
signed head counter retains absolute coordinates; the checked motion at
each nonhalting source boundary is zero. Table lookup, finite-word and list
codes, the absorbing step, and fixed-time iteration are primitive recursive.

A source frame records a source site, nine logical registers, a header gap
$g\ge0$, the first spare value $a_0$, and a finite spare tail $A$. Its tape
encoding $E_M(s)$ is a primitive-recursive function. The logical register
permutation in both delivered tables is
\[
(6,2,0,4,3,7,8,5,1).
\]
The retained independent verifier derives four list traversals directly
from the literal transition rows, without importing the compiler or using
an arithmetic-box oracle. Write
\[
B(v)=0\,1^{v+1},\qquad C(v)=1\,0\,1^v.
\]
The word identity $C(A)1=1B(A)$ follows by induction on the coded list
length. Its single-block identity and each self-loop crossing a run are
proved by quantifier-free counter induction. A list traversal has exact
duration $\operatorname{width}(A)=\sum_{v\in A}(v+2)$.
The checker composes these generic lemmas with literal transitions and
normalizes finite words. Nonnegative affine coefficients ensure that no
symbolic run of negative length is introduced. Its rules are mathematical
proof rules with the preceding primitive-recursive interpretation; the
Python implementation is not described as a proof-assistant kernel.

Every decrement has both its zero and successor branches. The source
contracts cover arbitrary natural registers, arbitrary first spare value,
arbitrary finite spare tail, and arbitrary nonnegative gap. The first spare
is present, so the complete spare file is nonempty. Address arithmetic and
interval-based tree resolution independently agree with the delivered
instruction graph. Every nonhalt result is checked to be the correct
padded target frame. Exact durations have thirteen slots: a constant, the
gap, nine register values, the first spare value, and the tail width.

The two starts are checked from an entirely blank tape, comparing the whole
finite set of one-cells. The source frame is at native state 84, head $-2$,
source site 25, and logical register vector $(2,0,\ldots,0)$.
Initially $g=0$, $a_0=1$, and $A$ consists of twenty-five ones; hence
$\operatorname{width}(A)=75$. The times are 57,870 and 57,712 for the two
new tables. The specified parents take 57,922 and 57,768 respectively.
Every transition through each new startup is also synchronized exactly
against its unaccelerated loop parent using the expanded clock above.

\subsection{Why source frames cover arbitrary physical halting times}
For one implementation let $s_j$ be the primitive-recursive iteration of
its augmented source step. Its logical projection is the common source
program. Let $h(s)$ mean that the source is ready to execute its halt
instruction; such a frame still has nonhalting physical control 84.
Let $k(s)\ge1$ be the exact normal-instruction duration, and $k_H(s)\ge1$
the halt-instruction duration, extended totally on irrelevant inputs.
With startup time $b$, define
\[
T(0)=b,\quad T(j+1)=T(j)+k(s_j),\qquad
P(j)=\prod_{i<j}(1-h(s_i)).
\]
The invariant is
\[
P(j)=1\ \Longrightarrow\
\bigl(\operatorname{Inv}(s_j)\land
 \operatorname{run}_M(T(j),\mathrm{blank})=E_M(s_j)
 \land T(j)\ge b+j\bigr).
\tag{*}
\]
All predicates here are Boolean combinations of primitive-recursive
functions. Quantifier-free induction proves (*) from the startup and the
complete source contracts. The factorization
$P(j+1)=P(j)(1-h(s_j))$ supplies the nonhalt premise needed for the next
contract. No quantification over all environments is hidden inside the
induction hypothesis.

If the logical source is ready to halt at index $n$, bounded first-hit
search finds the least such $m\le n$. Then $P(m)=1$, and the halt contract
provides a physical halt at $T(m)+k_H(s_m)$. Conversely suppose the physical
machine is halted by time $t$. If $P(t+1)=1$, (*) gives a nonhalting frame
at $T(t+1)\ge b+t+1>t$, contradicting absorption. Thus the finite Boolean
product $P(t+1)$ is zero. A bounded first-hit function returns a source
halt index $m\le t$. Its existence follows by quantifier-free induction
on a finite Boolean list. There is no use of Markov's principle or of
unbounded double-negation elimination. A hidden halt within a proposed
normal segment is excluded by its exact nonhalting endpoint under
absorbing semantics.

\subsection{An explicit polynomial witness bound}
The supplied coefficient audit checks, for every applicable contract of
the two new machines and their frozen parents, the conservative envelope
\[
3\le k(s),k_H(s)\le
1000+8g+4\sum_{i=0}^{8}r_i+2a_0+2\operatorname{width}(A).
\tag{D}
\]
The constant 1000 is deliberately larger than the measured maximum
coefficient. A positive-decrement branch uses $r_i-1$ in its variable
slot; replacing it by the actual input $r_i$ only increases the bound.

Each normal primitive increments at most one logical register by one.
The explicit frame updates give $g'=g+\delta+o$ with
$\delta,o\in\{0,1\}$. An increment appends one spare value 1; the other
primitives retain the spare list. After $j$ normal primitives, induction gives
\[
\sum_i r_i\le2+j,\qquad g\le2j,\qquad
 a_0=1,\qquad\operatorname{width}(A)\le75+3j.
\]
Substitution into (D) yields $k(s_j),k_H(s_j)\le1160+26j$.
Therefore $n$ normal primitives and their following halt instruction take
at most
\[
b+13(n+1)^2+1147(n+1).
\]
All four startup times are below 60,000. Define
\[
\boxed{\Theta(t)=60000+13(t+1)^2+1147(t+1)
                  =13t^2+1173t+61160.}
\]

\begin{theorem}[New relative iPRA result]
For each pairing
\[
(\mathrm{PA194LCG},\mathrm{PA194AX11GH}),\qquad
(\mathrm{PA194LCGS},\mathrm{PA194AX11GHS}),
\]
write the two machines as $N$ and $M$. Then iPRA proves
\[
\forall t\,(H_M(t)\to H_N(\Theta(t))),\qquad
\forall t\,(H_N(t)\to H_M(\Theta(t))).
\]
Consequently iPRA proves both their halting equivalence and their
nonhalting equivalence.
\end{theorem}
\begin{proof}
The source-array equality identifies the logical source iteration of
each pair. From a physical halting witness $t$ in either implementation,
the bounded extraction above gives a common logical halt index $n\le t$.
The other implementation halts within its displayed polynomial bound,
and hence by $\Theta(t)$. The construction and each induction are
primitive recursive and quantifier free. Existential introduction gives
the positive halting equivalence; ordinary intuitionistic reasoning gives
the negative equivalence. No premise about PA consistency occurs here.
\end{proof}
